% ============================================================================
\section{Metodología experimental: calibrar y validar}\label{sec:metodologia}
% ============================================================================

\subsection{Por qué hay que calibrar}

Las cinco metaheurísticas tienen parámetros (la temperatura inicial, la duración tabú, el
tamaño de la población\ldots) y su rendimiento depende mucho de ellos. Los apuntes no dan
valores porque no existen valores universales: dependen del problema, de las instancias y
del tiempo disponible. \textbf{Calibrar} es elegir esos valores mediante experimentos. Hay 21
parámetros en total y hay que justificar el valor de cada uno.

\subsection{Calibración y validación: qué instancias se usan para cada cosa}\label{sec:protocolo}

\begin{idea}[Respuesta corta]
Sí, es correcto: las instancias de la carpeta \cod{datos/} se usan para \textbf{calibrar} y
las de \cod{Validacion\_datos/} para \textbf{validar}. La regla que hace que esto funcione es
una sola: \emph{ninguna decisión se toma mirando las instancias de validación}.
\end{idea}

\paragraph{El problema que se quiere evitar.} Supóngase que se prueban 70 configuraciones
sobre las 30 instancias y se elige la mejor. Después se informa de lo bien que funciona esa
configuración\ldots{} sobre las mismas 30 instancias. Ese resultado es \emph{optimista}: la
configuración se eligió precisamente porque iba bien en esas instancias, y parte de su
ventaja puede deberse a la casualidad (a que encaja con las particularidades de esos
ficheros y de esas semillas) y no a que sea realmente mejor. Este fenómeno se llama
\textbf{sobreajuste}. Es el mismo problema que en estadística o en aprendizaje automático:
no se puede evaluar un modelo con los mismos datos con los que se ha ajustado.

\paragraph{La solución.} Separar las instancias en dos grupos que no se mezclan:
\begin{itemize}
  \item \textbf{Instancias de calibración} (también llamadas de \emph{entrenamiento} o de
        ajuste). Con ellas se prueban todas las configuraciones y se toman todas las
        decisiones.
  \item \textbf{Instancias de validación} (o de \emph{prueba}). Solo se usan al final, una
        vez fijados todos los parámetros, para medir cómo funciona la configuración elegida
        sobre instancias que \emph{no ha visto nunca}.
\end{itemize}
Es el procedimiento estándar en la configuración de algoritmos; así trabajan, por ejemplo,
las herramientas de configuración automática como irace~\cite{irace}, que distinguen
explícitamente entre instancias de entrenamiento e instancias de prueba.

\paragraph{Cómo se interpreta la validación.}
\begin{itemize}
  \item Si el rendimiento en validación es parecido al de calibración, los parámetros
        elegidos \emph{generalizan}: funcionan bien en instancias de este tipo, no solo en
        las seis con las que se eligieron.
  \item Si en validación es claramente peor, hay sobreajuste: la calibración ha aprendido
        particularidades de sus instancias.
  \item El resultado de validación es, además, la estimación honesta de la calidad de cada
        algoritmo. \textbf{Las cifras que deben ir al informe final son las de validación},
        no las de calibración.
\end{itemize}

\paragraph{Lo que no se puede hacer.} Mirar los resultados de validación y, a la vista de
ellos, volver a tocar un parámetro. En cuanto se hace eso, las instancias de validación han
participado en una decisión y dejan de ser independientes: se convierten, de hecho, en
instancias de calibración. Si se quisiera hacer una segunda ronda de ajuste haría falta un
tercer grupo de instancias reservado para la validación final. En esta práctica la
validación se ejecuta una sola vez, al terminar la calibración, y después no se modifica
nada.

\paragraph{El reparto 6 / 24.}

\begin{center}
\begin{tabular}{@{}llp{6.6cm}@{}}
\toprule
Carpeta & Instancias & Uso \\ \midrule
\cod{datos/} & SCP01, 06, 11, 16, 21, 26 (6) & Calibración: todas las decisiones. \\
\cod{Validacion\_datos/} & las 24 restantes & Validación: una sola ejecución final. \\
\bottomrule
\end{tabular}
\end{center}

\begin{decision}[Por qué pocas instancias para calibrar y muchas para validar]
\begin{itemize}[itemsep=1pt]
  \item \emph{Calibrar es caro}: cada configuración candidata se ejecuta sobre todas las
        instancias de calibración y con varias semillas. Con 6 instancias, 3 semillas y 5
        segundos, evaluar \emph{una} configuración son 90 segundos de cálculo, y hay que
        evaluar decenas.
  \item \emph{Validar es barato}: solo se ejecuta una configuración por algoritmo. Conviene
        hacerlo sobre muchas instancias para que la estimación sea fiable.
  \item \emph{Seis bastan porque las instancias son homogéneas}: las 30 tienen el mismo
        tamaño y la misma estructura (tabla~\ref{tab:instancias}). Si hubiera familias
        distintas (tamaños o densidades diferentes) habría que incluir representantes de
        cada una en ambos grupos.
  \item \emph{Las seis no se han escogido a dedo}: son una de cada cinco (01, 06, 11,
        \ldots), es decir, una muestra sistemática. Elegir ``las que mejor salen'' sería
        otra forma de sobreajuste.
\end{itemize}
\end{decision}

El mismo criterio se ha aplicado a las decisiones de los algoritmos básicos: el criterio
voraz (sección~\ref{sec:res-voraz}) y el número de entornos de la búsqueda local se han
elegido con las 6 instancias de calibración.

\subsection{Medida de calidad: RPD}\label{sec:rpd}

El coste óptimo no es el mismo en todas las instancias (va de 322 a 391), así que los costes
no se pueden promediar ni comparar directamente. Se usa la \textbf{desviación porcentual
relativa} respecto del óptimo (RPD, \emph{Relative Percentage Deviation}):
\[
  \rpd \;=\; 100\cdot\frac{z-z^\ast}{z^\ast}\ \%,
\]
donde $z$ es el coste obtenido y $z^\ast$ el óptimo de la instancia
(sección~\ref{sec:optimos}). Un RPD de 0 significa que se ha encontrado el óptimo; un RPD de 4
significa un 4\,\% más caro que el óptimo. Menor es mejor. El RPD de una configuración es la
\emph{media} de los RPD de todas sus ejecuciones.

\subsection{Semillas}

Las metaheurísticas usan números aleatorios, de modo que dos ejecuciones dan resultados
distintos. Juzgar una configuración por una sola ejecución sería juzgarla en parte por
suerte. Cada configuración se ejecuta con 3 semillas (1, 2 y 3) en cada instancia:
\begin{itemize}
  \item calibración: $6\times 3=18$ ejecuciones por configuración;
  \item validación: $24\times 3=72$ ejecuciones por algoritmo.
\end{itemize}
Las semillas son las mismas para todas las configuraciones, de modo que la comparación entre
ellas es lo más pareja posible.

\subsection{Presupuesto de cálculo}

Todas las ejecuciones de todas las metaheurísticas tienen el mismo límite: \textbf{5 segundos
de CPU}. Darles el mismo tiempo es lo que hace justa la comparación entre algoritmos.

\begin{idea}[Advertencia: los parámetros dependen del tiempo]
Los valores calibrados son los mejores \emph{para 5 segundos}. Algunos parámetros (sobre
todo $\alpha$ y $L$ del temple, que fijan a qué velocidad se enfría) están directamente
ligados al tiempo disponible. Si el informe final usara otro límite de tiempo, habría que
recalibrarlos.
\end{idea}

Las ejecuciones de referencia de esta memoria se han hecho en un procesador Intel Xeon a
2.1~GHz, con el programa compilado con \cod{gcc -O2}. En un ordenador más rápido el programa
hace más iteraciones en los mismos 5 segundos y los resultados son algo mejores; por eso, si
se repiten los experimentos en otra máquina, hay que repetir \emph{todos} (calibración y
validación) en la misma y regenerar las tablas (sección~\ref{sec:uso}).

\subsection{Procedimiento: un factor cada vez}\label{sec:ofat}

\paragraph{Por qué no se prueban todas las combinaciones.} El temple tiene 5 parámetros. Con
3 o 4 valores para cada uno hay $4\cdot4\cdot4\cdot3\cdot3=576$ combinaciones; a 90 segundos
cada una, más de 14 horas solo para el temple. No es viable.

\paragraph{El procedimiento.} Para cada algoritmo:
\begin{enumerate}
  \item Se parte de una \textbf{configuración base} razonable y de una lista de valores
        candidatos para cada parámetro.
  \item Se toma el primer parámetro y se prueban todos sus valores candidatos,
        \emph{dejando los demás parámetros como están}. Se fija el valor con menor RPD
        medio.
  \item Se pasa al segundo parámetro, y así hasta el último. Cada paso se llama
        \emph{etapa}.
\end{enumerate}
En caso de empate se conserva el valor que ya estaba. Este procedimiento se conoce como
\emph{un factor cada vez} (OFAT, \emph{one factor at a time}). Con él, el temple necesita
evaluar 14 configuraciones en lugar de 576, y las cinco metaheurísticas juntas,
56%
 configuraciones (1008%
 ejecuciones).

\begin{decision}[Por qué este procedimiento]
Es el más barato que permite justificar \emph{cada} parámetro por separado con una tabla: se
ve qué valores se probaron, qué resultado dio cada uno y por qué se eligió el que se eligió.
Además muestra la \emph{sensibilidad}: si al mover un parámetro el RPD apenas cambia, ese
parámetro importa poco; si cambia mucho, es crítico.
\end{decision}

\paragraph{Valores candidatos.} Para cada parámetro se han tomado 3 o 4 valores que cubren
de ``poco'' a ``mucho'', incluyendo cuando existe el valor que \emph{desactiva} el mecanismo
($\beta=0$: sin memoria a largo plazo; $\lambda_f=1$: sin memoria en GRASP; \cod{pmut}~$=0$:
sin mutación; \cod{bl}~$=0$: sin mejora local; \cod{reinicio} enorme: sin intensificación).
Así la calibración sirve también para comprobar si cada componente del algoritmo aporta algo.

\paragraph{Cuándo una diferencia es real.} Dos ejecuciones de la misma configuración sobre la
misma instancia con distinta semilla difieren, de media, en una o dos décimas de RPD
(desviación típica entre semillas de 0.08 a 0.27 según el algoritmo). Al promediar 18
ejecuciones ese ruido se reduce, pero no desaparece: \textbf{diferencias de menos de unas
cinco centésimas entre dos valores de un parámetro no son significativas} y deben leerse como
un empate.

\paragraph{Limitaciones, y la segunda vuelta opcional.} El método tiene dos debilidades
conocidas:
\begin{itemize}
  \item Ignora las \emph{interacciones}: el mejor valor de un parámetro puede depender del
        valor de otro, y los primeros parámetros se eligen cuando los últimos aún están en
        su valor base.
  \item Solo explora los valores de la lista. Si gana el valor más grande (o el más
        pequeño) de la lista, puede que el óptimo esté más allá.
\end{itemize}
Para atenuarlas, el guion \cod{calibrar.py} admite una \textbf{segunda vuelta} opcional
(\cod{VUELTAS=2 python3 calibrar.py}): repasa otra vez todos los parámetros con los demás ya
en sus valores elegidos y, cuando el valor vigente es un extremo de su lista, añade un valor
más en esa dirección. Duplica aproximadamente el tiempo de cálculo. Los resultados de esta
memoria corresponden a \textbf{una vuelta}; en la sección~\ref{sec:calibracion} se señala en
qué parámetros la segunda vuelta tendría más interés. Para un estudio más ambicioso, la
alternativa actual es la configuración automática con herramientas como
irace~\cite{irace}, que reparten el esfuerzo de cálculo de forma estadísticamente eficiente.

\subsection{Reproducibilidad}

Todo el experimento se puede repetir con tres órdenes (sección~\ref{sec:uso}). Los guiones
guardan cada ejecución en un fichero CSV en cuanto termina y, si se interrumpen, continúan
donde lo dejaron sin repetir nada. Las tablas y figuras de esta memoria las genera
\cod{generar\_tablas.py} a partir de esos ficheros: ninguna cifra de las tablas se ha copiado
a mano.
